\documentclass{beamer}
\usepackage{tikz,tikz-cd,amsmath,amssymb,float,graphicx}
\usepackage{quiver}
\usepackage{XCharter}\usepackage{beramono}

\let\d\relax
\DeclareMathOperator{\Hom}{Hom}
\DeclareMathOperator{\Sub}{SubPol}
\DeclareMathOperator{\Set}{\textbf{Set}}
\DeclareMathOperator{\Pp}{\mathcal{P}}
\DeclareMathOperator{\Ll}{\mathcal{L}}
\DeclareMathOperator{\Cc}{\mathcal{C}}
\DeclareMathOperator{\Q}{\mathbb{Q}}
\DeclareMathOperator{\R}{\mathbb{R}}
\DeclareMathOperator{\Z}{\mathbb{Z}}
\DeclareMathOperator{\d}{\lozenge}
\newcommand{\upred}{
  \mathrel{\circ}~\kern-0.6em~\relbar\joinrel\twoheadrightarrow
}
\newcommand{\fupred}{
  \mathrel{\circ}~\kern-0.6em~\relbar\joinrel\mapsto
}
\newcommand{\set}[1]{\left\{#1\right\}}
\usetheme[progressbar=frametitle]{metropolis}

\title{A Goldblatt-Thomason Theorem for Compact Polyhedra}
\subtitle{}
\author{David Gabelaia, Francesco Tognetti}
\date{July 2026}
\institute{TACL @ Krakow}



\begin{document}
\begin{frame}
    \maketitle
\end{frame}
\begin{frame}{Hello!}
    A small presentation, \\
    from least to most important 
    \begin{itemize}
        \item My name is \textbf{Francesco Tognetti},
        \item I'm a second year PhD student 
        at the \textbf{University of Luxembourg},
        \item I have a very cute \textbf{dog}.
    \end{itemize}
\end{frame}
\begin{frame}{Oopsie}
    The \textbf{abstract} in the conference website 
    said I'd be presenting a Goldblatt-Thomason theorem for 
    \textbf{Pre-compact} polyhedra.

    Apparently that approach was \textbf{flawed}, to the point that
    the result turned out not to be valid. 

    We managed to prove one for \textbf{Compact} polyhedra instead.
    \vfill
    (I'm very sorry for the mix-up)
\end{frame}
\section{Topological and Polyhedral semantics}
\begin{frame}{Topological semantics}
    A quick recap is in order:

    We are doing \textbf{Polyhedral semantics of modal logic}, that can be 
    seen as a slight modification of \textbf{Topological semantics}.

    We have the standard unimodal language 
    \[\Ll=\set{p\in P\ \|\ \bot\ \|\ \lnot\varphi\ \|\ \varphi\vee\psi\ \|\ \d\varphi}\]
\end{frame}
\begin{frame}{Polyhedral Semantics}
    As it's the case in topological semantics we interpret formulas in a space $X$
    \[\nu:\Ll\to \Pp(X)\]
    \begin{itemize}
        \item $\bot$ as the \textbf{Empty Set}\hfill $\nu(\bot)=\emptyset$
        \item $\lnot$ as the \textbf{Complementary } \hfill $\nu(\lnot\varphi)=X\setminus \nu(\varphi)$
        \item $\vee$ as the \textbf{Union} \hfill $\nu(\varphi\vee\psi)=\nu(\varphi)\cup\nu(\psi)$
        \item $\d$ as the \textbf{Topological closure} \hfill $\nu(\d\varphi)=\overline{\nu(\varphi)}$
    \end{itemize}
\end{frame}
\begin{frame}{Polyhedral Semantics}
    The distinguishing factor between \textbf{Topological} and \textbf{Polyhedral} semantics
    is that
    \begin{itemize}
        \item We require $X$ to be a \textbf{Polyhedron}
        \item For any $\varphi$ we require $\nu(\varphi)$ to be a \textbf{Subpolyhedron}
    \end{itemize}
\end{frame}
\begin{frame}{Polyhedra}
    Polyhedra are essentially higher dimensional generalizations of polygons,
    thus the nicest way to describe them is as \textbf{unions of simplices}.
    \begin{itemize}
        \item A $0-$Simplex is a point,
        \item A $1-$Simplex is a segment,
        \item A $2-$Simplex is a triangle,
        \item A $3-$Simplex is a tetrahedron,
        \item $\cdots$
    \end{itemize}
\end{frame}
\begin{frame}{Polyhedra}
    Since we want to allow \textbf{open} sets (in the euclidean topology)
    What we consider as ``acceptable subsets'' is 
    \textbf{unions of relative interior of simplices}
    \vfill
\end{frame}
\begin{frame}{Compactness}
    It's not obvious what an \textbf{infinite}-polyhedron is:
    If we allowed them to be arbitrary unions of (relative interiors of) simplices then 
    any topological space is a polyhedron, as every point is a simplex.

    There are many possible generalisations but
    
    for now we may restrict to the \textbf{compact} case.
\end{frame}
\begin{frame}
    This gives rise to a much \textbf{tamer} topology:

    For example every polyhedron of dimension $\leq n$ 
    validates $Bd_n$ (the bounded depth formula of order $n$)
    \vfill
\end{frame}
\begin{frame}
    Or any convex polyhedron validates $\chi(\ \ \ \ \ \ )+\chi(\ \ \ \ \ \ )$
    \vfill 
\end{frame}
\section{A Definability theorem}
\begin{frame}{Definability}
    From here the \textbf{definability} question is pretty natural:
    Which \textbf{conditions} are necessary and sufficient for a class of polyhedra to be 
    \textbf{modally definable}?
    \begin{theorem}[GTKP]
        Let $\Cc$ be a class of compact polyhedra. Then $\Cc$ is modally definable
        if and only if it's closed under 
        \begin{enumerate}[\textbf{(1)}]
            \item \textbf{Finite disjoint unions},
            \item \textbf{Face up-reductions}
        \end{enumerate}
    \end{theorem}
\end{frame}
\begin{frame}{Finite disjoint unions}
    \begin{definition}
        Let $X, Y$ be two sets, we define the 
        \textbf{Disjoint union} of $X$ and $Y$ as the colimit diagram 
        \[\begin{tikzcd}[ampersand replacement=\&]
	\& U \& \\
	X \& {X+Y} \& Y
	\arrow[from=2-1, to=1-2]
	\arrow[from=2-1, to=2-2]
	\arrow["{\exists!}"{description}, dashed, from=2-2, to=1-2]
	\arrow[from=2-3, to=1-2]
	\arrow[from=2-3, to=2-2]
\end{tikzcd}\]
        In the category of sets
    \end{definition}
    Blah blah words to fill out the page 
    because otherwise the visual impact doesn't reflect the
    joke appropriately du dudududu du du
    I hear the drums echoing tonight 
\end{frame}
\begin{frame}{ahah}
    \begin{center}
        Just kidding, I know you know
    \end{center}
\end{frame}
\begin{frame}{Face up-reductions}
    To understand what a face up-reduction is, we start from the concept 
    of topological up-reduction:

    \begin{definition}[Up-reduction]
        Let $X,Y$ be topological spaces, an \textbf{up-reduction} $f:X\upred Y$
        Is a \textbf{Surjective partial map} such that 
        \begin{itemize}
            \item The \textbf{domain} of $f$ is an open subset $U\subseteq^\circ X$,
            \item $f_{|U}$ is an \textbf{open continuous map}.
        \end{itemize}
    \end{definition}
\end{frame}
\begin{frame}{Face posets}
    The reason why polyhedra are so \textbf{well-behaved} is that we can 
    \textbf{triangulate} them properly.

    That is, for any polyhedron $P\subseteq \R^n$ there are (uncountably) \textbf{many} ways
    of choosing a \textbf{triangulation}, i.e. a collection of simplices 
    that cover the polyhedron.
    \vfill
\end{frame}
\begin{frame}
    And furthermore, chosen a finite set of subpolyhedra $P_1,...,P_k$ 
    there always exists a finite triangulation of $P$ that triangulates all of them 

    \vfill 
\end{frame}
\begin{frame}{Finiteness}
    This means that despite the algebra of subpolyhedra $\Sub(P)$ being infinite,

    given $\varphi_1,...,\varphi_n$, we can always find a \textbf{finite} subalgebra 
    that contains $\set{\nu(\varphi_1),...,\nu(\varphi_n)}$ 
\end{frame}
\begin{frame}{Posets}
    Triangulations have a natural poset structure, given by the \textbf{``face of''} relation:
    \vfill 
\end{frame}
\begin{frame}
    For example
    \vfill 
\end{frame}
\begin{frame}{Poset topology}
    Every poset comes equipped with a topology, where 
    the \textbf{upward-closed} sets are open, and the 
    \textbf{downward-closed} sets are closed.
    \vfill 
\end{frame}
\begin{frame}
    This is compatible with our topology on polyhedra: 
    Given a triangulation $T$ of a polyhedron $P$, (euclidean) 
    \textbf{open subpolyhedra} form upward-closed sets in $T$
    \vfill
\end{frame}
\begin{frame}{Face up-reductions}
    Given two polyhedra $P$ and $Q$ a \textbf{face up-reduction}
    $P\fupred Q$
    
    Is a surjective partial map $P\twoheadrightarrow \Sigma$ to a triangulation of $Q$
    such that the \textbf{preimage} of every open of $\Sigma$ is an 
    open subpolyhedron of $P$.
    \vfill 
\end{frame}
\begin{frame}{In particular}
    Given any polyhedron $Q$ and any triangulation $\Sigma$, then 
    \[Q\fupred \Sigma\]
    \vfill 
\end{frame}
\section{Proof}
\begin{frame}{The proof, essentially}
    \textbf{Modally definable $\implies$ closed under (1) and (2):}
    Disjoint unions are trivial, let's focus on face up-reductions;

    Suppose $P\models F$ and $P\fupred Q$ through a triangulation $\Sigma$, then assume $Q\not{\models}F$
    
    Then there exists a formula $\varphi\in F$ such that $Q\not\models \varphi$

    But then, we can find a point $q$ of $Q$ such that $Q,q\not\models\varphi$
    \vfill 
\end{frame}
\begin{frame}
    This means that there exists a triangulation $\Sigma'$ that triangulates $\nu(\varphi)$
    and subdivides $\Sigma$:
    We can extend $P\upred\Sigma$ to $P\upred\Sigma'$
    \vfill
    \vfill
    Meaning $P$ refutes $\varphi$ as well\hfill $\to\gets$
\end{frame}
\begin{frame}
    \textbf{Closed under (1) and (2) $\implies$ modally definable:}

    Let $\Cc$ be a class closed under \textbf{(1)} and \textbf{(2)}.

    Call $\Ll_{\Cc}$ the logic of $\Cc$ and call $\Cc'$ the polyhedra validating $\Ll_{\Cc}$.

    \vfill 

    We only need to show that $\Cc'\subseteq \Cc$, so let $Q\in \Cc'$.
\end{frame}
\begin{frame}
    $Q\models \Ll$.  We want to construct $P\in \Cc$ such that $P\fupred Q$.

    To do that, we can pick a triangulation $\Sigma$ of $Q$, and consider all of its rooted
    open subposets $(\uparrow v_i)$.
    \vfill
    We can pick the Jankov-Fine formulas $\chi_i=\chi(\uparrow v_i)$for each of them, \textbf{excluding
    the presence of an up-reduction to $(\uparrow v_i)$}.
\end{frame}
\begin{frame}
    As $(\uparrow v_i)$ are generated subframes of $\Sigma$, $\Sigma\not\models\chi_i$
    Thus none of the $\chi_i$ belong to $\Ll$

    We can find polyhedra $P_i\in\Cc$ up-reducing to each $(\uparrow v_i)$
    \vfill
    Thus $\bigsqcup P_i\upred \bigsqcup(\uparrow v_i)\upred \Sigma$
    \vfill
    Thus $P:=\bigsqcup P_i\fupred Q$
\end{frame}
\section{Boring part over}
\begin{frame}{What can we do with this?}
    As expected, we can show that the class of polyhedra of bounded dimension is 
    modally definable:
    \[\dim(P)\leq n\iff P\models Bd_n\]
\end{frame}
\begin{frame}
    Less expectedly, the class of convex polyhedra is \textbf{not} modally definable:
    There must exist compact non-convex polyhedra that validate the same formulas as 
    all convex ones.
\end{frame}
\begin{frame}
    Even more unexpectedly, we can talk about some form of homology in one dimension:
    
    The class of 1d-polyhedra with less than $n$ holes for connected component is modally
    definable
\end{frame}
\begin{frame}
    Somewhat more expectedly, the class of (locally) contractible polyhedra is not modally
    definable.
\end{frame}
\section{future work}
\begin{frame}{up-reductions}
    There is \textbf{some hope} that we may get a result closer to the topological one:

    Every \textbf{up-reduction} implies \textbf{face up-reductions}
    \vfill 
    And the existence of a \textbf{face up-reduction} is independent on the choice of 
    face poset for the same polyhedron.
    \vfill
    The converse implication is still unknown 
\end{frame}
\begin{frame}{classes}
    It would be interesting to see if applying conditions that generalize the 
    \textbf{``less than n holes''} in higher dimensions are possible.
\end{frame}
\section{Thank you for listening!}
\end{document}